\begin{proof}[推论~\ref{cor:quantitative-circuits} 的证明]
我们使用 Williams 的见证电路论证~\cite[定理~6.1]{Wil13}，
并结合 Jahanjou、Miles 和
Viola 的局部归约~\cite[定理~1 和第~6 节]{JMV13}。

根据 \v{Z}\'ak 的非确定性时间层次定理~\cite{Zak83}，采用
~\cite[完整版中的定理~2.1]{Wil13} 所陈述的形式，固定一个一元语言
$H\in\operatorname{NTIME}(2^n)\setminus\operatorname{NTIME}(2^{n/2})$，
它定义在多带图灵机上。对所有充分大的 $n$，将局部归约应用于其固定机器在 $1^n$ 上的运行，
局部性参数取 $r=2$。时间配额
$T=2^n n^{O(1)}$ 包含了在归约中模拟该机器的多对数级开销，
并给出 $q=n+O(\log n)$ 个子句索引位。该归约在多项式时间内构造一个解码器，
它在每个索引上给出三个变量标签及其符号。令 $\varphi_n$
由全部 $2^q$ 个索引上解码得到的子句组成。它可满足
当且仅当 $1^n\in H$。

解码器的 $O(n)$ 个输出位中的每一个都具有常数深度决策
树，因此至多依赖于固定数目 $w$ 个索引位。
在所有汉明重量至多 $w$ 的串上对该解码器求值
即可恢复每个支撑集。事实上，见证对
某一坐标的依赖的两个串可以在支撑集之外置零，同时保留
它们不同的输出。同样的求值给出每个支撑集上的
真值表。这在多项式时间内构造出一个由 $O(n)$ 个二输入门
构成的解码器。

对固定的正整数 $d$，将每个变量标签填充到 $dn$ 位，
并声明全部 $2^{dn}$ 个变量。若对某个这样的 $n$ 有 $|i|=dn$ 且
$\varphi_n$ 可满足，则令 $W(i)$ 为其
字典序第一个满足赋值的第 $i$ 位，其中变量按其
标签排序。在所有其他情形令 $W(i)=0$。$i$ 的
长度确定 $n$。为判定 $W$，写出这 $2^q$ 个子句，向
NP 谕示询问该公式是否可满足；若可满足，
则通过询问其当前前缀后接零是否具有可满足的扩展
来恢复第一个赋值。至多 $2^{dn}$ 次长度为 $2^{O(n)}$ 的前缀
询问即足够，故 $W\in\mathrm E^{\mathrm{NP}}$。

令 $a(t)=\sqrt{\log_2 t}$。检验定理~\ref{thm:main} 的
常数密度部分中的计数是否为正，可在时间
$V^b2^{A V/a(V)}$ 内判定具有 $V$ 个变量、至多 $5V$ 个子句的 $3$-SAT，其中 $A,b>0$ 固定。我们将
对数底数的变化吸收进 $A$。该求解器使用多项式空间，因此
多带模拟可以在每次位访问时以代价 $V^{O(1)}$ 扫描其工作空间，
该代价被吸收进 $b$。

固定正有理数 $\eta<1/(12A)$，并对 $n\ge2$ 令
$s_n=\lfloor\eta n\lfloor a(n)\rfloor\rfloor$。
该阈值可利用 $n$ 的二进制长度和
$\lfloor\log_2 n\rfloor$ 的整数平方根计算。假设对
所有充分大的 $n$ 有 $C_W(dn)\le s_n$。为判定 $H$，
拒绝非一元输入。在 $1^n$ 上，猜测一个作用于 $dn$ 个索引位、
至多 $s_n$ 个门的电路。将解码器的三个变量标签
输入该电路的三个副本，并检验被索引的子句是否
被违反。若测试电路不可满足则接受。当
$1^n\in H$ 时，计算长度为 $dn$ 的 $W$ 的电路给出一个
可接受的猜测。当 $1^n\notin H$ 时，每个猜测都违反某个子句。

测试电路有 $q$ 个输入和至多 $3s_n+O(n)$ 个门。上述
编码为每个输入和门使用一个变量，为每个门至多使用四个
宽度至多三的子句以强制其真值表，
并使用一个单位子句固定输出。用未使用变量
填充后，对固定整数 $B$ 得到 $U_n=3s_n+Bn$ 个变量和至多
$5U_n$ 个子句。由于 $U_n=(3\eta+o(1))n a(n)$ 且 $a(U_n)\sim a(n)$，
该检验耗时 $2^{(3A\eta+o(1))n}$。猜测和构造
电路花费多项式时间，因此对所有充分大的 $n$，$H$ 可在时间
$2^{n/2}$ 内非确定性判定。有限多个
例外长度可以被硬连线，这与 $H$ 的选取矛盾。
因此对无穷多个 $n$ 有 $C_W(dn)>s_n$。由于
$s_n\sim\eta n\sqrt{\log_2 n}$ 且 $d$ 固定，取 $L=W$ 及任意
$0<c<\eta/(d\sqrt{\log 2})$ 便证明了 \eqref{eq:quantitative-circuit-lower-bound}。
\end{proof}
